% Canonical PDF foundation for books in the NeutrinoHit universe.

\usepackage{fontspec}
\usepackage{polyglossia}
\usepackage{unicode-math}
\usepackage{amsmath}
\usepackage{mathtools}
\usepackage{graphicx}
\usepackage{float}
\usepackage{booktabs}
\usepackage{longtable}
\usepackage{array}
\usepackage{microtype}
\usepackage{hyperref}
\usepackage{tikz}
\usepackage[most]{tcolorbox}
\usepackage{caption}
\usepackage{ragged2e}
\usepackage{etoc}
\usepackage[automark]{scrlayer-scrpage}

\defaultfontfeatures{Ligatures=TeX, Scale=MatchLowercase}
\setmainfont{Libertinus Serif}
\setsansfont{Libertinus Sans}
\setmonofont{Libertinus Mono}[Scale=0.9]
\setmathfont{Libertinus Math}

\definecolor{BookInk}{HTML}{18262F}
\definecolor{BookInkSoft}{HTML}{42545F}
\definecolor{BookMuted}{HTML}{687984}
\definecolor{BookPaper}{HTML}{FFFFFF}
\definecolor{BookCanvas}{HTML}{F4F6F7}
\definecolor{BookLine}{HTML}{D9E0E3}
\definecolor{BookLineStrong}{HTML}{BCC9CE}
\definecolor{BookTeal}{HTML}{0F6972}
\definecolor{BookTealDark}{HTML}{094B52}
\definecolor{BookTealSoft}{HTML}{E8F1F2}
\definecolor{BookBlue}{HTML}{2F6F9F}
\definecolor{BookBlueSoft}{HTML}{EDF3FA}
\definecolor{BookCoral}{HTML}{C7472F}
\definecolor{BookCoralSoft}{HTML}{FAEEEB}
\definecolor{BookGold}{HTML}{B97817}
\definecolor{BookGoldSoft}{HTML}{FBF3E4}
\definecolor{BookBiography}{HTML}{465B8C}
\definecolor{BookBiographySoft}{HTML}{F0F2F8}
\definecolor{BookMargin}{HTML}{F1F3F4}

% Compatibility aliases let established books keep their local command style.
\colorlet{bookInk}{BookInk}
\colorlet{bookMuted}{BookMuted}
\colorlet{bookTeal}{BookTeal}
\colorlet{bookTealSoft}{BookTealSoft}
\colorlet{bookBlue}{BookBlue}
\colorlet{bookBlueSoft}{BookBlueSoft}
\colorlet{bookCoral}{BookCoral}
\colorlet{bookCoralSoft}{BookCoralSoft}
\colorlet{bookBiography}{BookBiography}
\colorlet{bookBiographySoft}{BookBiographySoft}
\colorlet{bookMargin}{BookMargin}

\setlength{\emergencystretch}{3em}
\raggedbottom

\hypersetup{
  pdfborder={0 0 0},
  colorlinks=true,
  linkcolor=BookTeal,
  citecolor=BookTeal,
  urlcolor=BookTeal
}

% Shared semantic boxes used by biographies, media cards, notes and exercises.
\newtcolorbox{semanticpdfbox}[1][]{
  enhanced,
  breakable,
  sharp corners,
  colback=BookTealSoft,
  colframe=BookTeal,
  boxrule=0pt,
  borderline west={1.2mm}{0pt}{BookTeal},
  left=4mm,
  right=3mm,
  top=2mm,
  bottom=2mm,
  #1
}

% Public labels and colour hook. A book changes these without copying the
% layout machinery: for example, "Лекция" in the neutrino book and "Глава"
% in the statistics book.
\providecommand{\bookchaptercolor}{BookTeal}
\newcommand{\BookChapterLabel}{Глава}
\newcommand{\BookPartLabel}{РАЗДЕЛ КНИГИ}
\newcommand{\BookPartContentsLabel}{Главы}
\newcommand{\BookTocChapterLabel}{Глава}

\clearpairofpagestyles
\ihead{\headmark}
\newcommand{\BookHeaderPage}{\ifnum\value{chapter}=0\pagemark\fi}
\ohead{\BookHeaderPage}
\setkomafont{pageheadfoot}{\small\sffamily\color{BookMuted}}
\setkomafont{pagehead}{\small\sffamily\color{BookInk}}
\KOMAoptions{headsepline=0.35pt}
\addtokomafont{chapter}{\sffamily\color{BookInk}}
\addtokomafont{section}{\sffamily\color{BookInk}}
\addtokomafont{subsection}{\sffamily\color{BookInk}}

% Numbered chapter opener with a local table of contents.
\renewcommand*{\chapterformat}{}
\renewcommand*{\chapterlinesformat}[3]{%
  \ifnum\value{chapter}>0\relax
    \parbox[t]{\linewidth}{%
      {\sffamily\bfseries\color{\bookchaptercolor}%
       \fontsize{15}{18}\selectfont \BookChapterLabel~\thechapter\par}%
      \vspace{7mm}%
      {\sffamily\bfseries\color{BookInk}%
       \fontsize{34}{38}\selectfont #3\par}%
      \vspace{9mm}%
      {\color{\bookchaptercolor}\rule{\linewidth}{1.3mm}}%
    }%
  \else
    \parbox[t]{\linewidth}{#3}%
  \fi
}
\RedeclareSectionCommand[beforeskip=10mm,afterskip=0pt]{chapter}

% Part opener with the list of chapters belonging to the current part.
\renewcommand*{\partformat}{}
\renewcommand*{\partlineswithprefixformat}[3]{%
  \parbox[t]{\linewidth}{%
    {\sffamily\bfseries\color{\bookchaptercolor}%
     \fontsize{13}{16}\selectfont \BookPartLabel\par}%
    \vspace{9mm}%
    {\sffamily\bfseries\color{BookInk}%
     \fontsize{32}{37}\selectfont #3\par}%
    \vspace{11mm}%
    {\color{\bookchaptercolor}\rule{\linewidth}{1.5mm}}\par
    \vspace{13mm}%
    {\sffamily\bfseries\color{BookInk}%
     \fontsize{18}{22}\selectfont \BookPartContentsLabel\par}%
    \vspace{3mm}%
    {\color{\bookchaptercolor}\rule{18mm}{0.8mm}}\par
    \vspace{5mm}%
    \begingroup
      \normalfont
      \setlength{\parindent}{0pt}%
      \etocsettocstyle{}{}%
      \etocsetnexttocdepth{chapter}%
      \localtableofcontents
    \endgroup
  }%
}
\RedeclareSectionCommand[beforeskip=20mm,afterskip=0pt]{part}

\newcommand{\BookHidePartNumber}[1]{}
\newcommand{\BookChapterTocNumber}[1]{\BookTocChapterLabel~#1}
\DeclareTOCStyleEntry[numwidth=0pt,entrynumberformat=\BookHidePartNumber]{tocline}{part}
\DeclareTOCStyleEntry[numwidth=5.6em,entrynumberformat=\BookChapterTocNumber]{tocline}{chapter}

\newcommand{\bookchapteropeningcontents}{%
  \vspace{14mm}%
  \noindent
  {\sffamily\bfseries\color{BookInk}\fontsize{18}{22}\selectfont
   Содержание\par}%
  \vspace{3mm}%
  {\color{\bookchaptercolor}\rule{18mm}{0.8mm}}%
  \vspace{5mm}%
  \begingroup
    \setlength{\parindent}{0pt}%
    \etocsettocstyle{}{}%
    \etocsetnexttocdepth{section}%
    \localtableofcontents
  \endgroup
  \vfill
  \clearpage
}

% Outer navigation rail and circular page number. Dimensions are shared; the
% colour follows the current part through \bookchaptercolor.
\newif\ifbookpagenavigator
\bookpagenavigatortrue
\newlength{\bookpagerailwidth}
\newlength{\bookpagerailcenter}
\newlength{\bookpagetabheight}
\newlength{\bookpagetabcenter}
\newlength{\bookpagenumberdiameter}
\setlength{\bookpagerailwidth}{48mm}
\setlength{\bookpagerailcenter}{24mm}
\setlength{\bookpagetabheight}{28mm}
\setlength{\bookpagetabcenter}{14mm}
\setlength{\bookpagenumberdiameter}{17mm}
\newcommand{\bookpagerailcolor}{BookMargin}
\newcommand{\bookpagetabcolor}{\bookchaptercolor}
\newcommand{\bookpagenumberfill}{white}
\newcommand{\bookpagenumbercolor}{\bookchaptercolor}

\tikzset{
  book page rail/.style={fill=\bookpagerailcolor,draw=none},
  book page tab/.style={fill=\bookpagetabcolor,draw=none},
  book page number/.style={
    circle,
    fill=\bookpagenumberfill,
    draw=none,
    minimum size=\bookpagenumberdiameter,
    inner sep=0pt,
    font=\sffamily\bfseries\fontsize{22}{22}\selectfont,
    text=\bookpagenumbercolor
  }
}

\newcommand{\BookPageNavigatorLeft}{%
  \path[book page rail]
    (current page.north west) rectangle
    ([xshift=\bookpagerailwidth]current page.south west);
  \path[book page tab]
    (current page.north west) rectangle
    ([xshift=\bookpagerailwidth,yshift=-\bookpagetabheight]current page.north west);
  \node[book page number]
    at ([xshift=\bookpagerailcenter,yshift=-\bookpagetabcenter]current page.north west)
    {\thepage};
}

\newcommand{\BookPageNavigatorRight}{%
  \path[book page rail]
    ([xshift=-\bookpagerailwidth]current page.north east) rectangle
    (current page.south east);
  \path[book page tab]
    ([xshift=-\bookpagerailwidth]current page.north east) rectangle
    ([yshift=-\bookpagetabheight]current page.north east);
  \node[book page number]
    at ([xshift=-\bookpagerailcenter,yshift=-\bookpagetabcenter]current page.north east)
    {\thepage};
}

\newcommand{\BookRenderPageNavigator}{%
  \ifbookpagenavigator
    \ifnum\value{chapter}>0\relax
      \begin{tikzpicture}[remember picture,overlay]
        \ifodd\value{page}
          \BookPageNavigatorRight
        \else
          \BookPageNavigatorLeft
        \fi
      \end{tikzpicture}%
    \fi
  \fi
}
\AddToHook{shipout/background}{\BookRenderPageNavigator}
